% Working draft: does not alter theorem_DK.tex or any certification flag.
\documentclass[11pt]{amsart}
\usepackage{amsmath,amssymb,amsthm,mathtools,enumitem,hyperref}
\newtheorem{theorem}{Theorem}[section]
\newtheorem{proposition}[theorem]{Proposition}
\newtheorem{lemma}[theorem]{Lemma}
\newtheorem{corollary}[theorem]{Corollary}
\theoremstyle{definition}
\newtheorem{assumption}[theorem]{Assumption}
\newtheorem{remark}[theorem]{Remark}
\newcommand{\HH}{\mathbb H^3}
\newcommand{\CC}{\mathbb C}
\newcommand{\ZZ}{\mathbb Z}
\newcommand{\G}{\Gamma}
\newcommand{\Hsub}{\Gamma_0(2+i)}
\newcommand{\norm}[1]{\lVert #1\rVert}
\newcommand{\abs}[1]{\lvert #1\rvert}
\newcommand{\rhoind}{\rho_{\mathrm{ind}}}
\newcommand{\D}{\Delta}
\title{A six-copy/two-cusp defect theorem for $\Gamma_0(2+i)$}
\author{}
\date{Working draft, 2026-07-13}
\begin{document}
\begin{abstract}
The corrected six-copy field is a rank-six induced permutation local
system on the level-one Picard orbifold. We prove the finite-index
identification and the two singular cusp channels, and formulate a
quantitative smooth-gluing theorem whose inputs can be interval certified.
The continuous spectrum is removed either by an exact oldspace,
quarter-turn-odd projector for the present target or, in general, by a
separately normalized Hecke cuspidalizer. The existing value-only defect is
not yet an admissible theorem input.
\end{abstract}
\maketitle

\section{Why the scalar proof cannot merely be multiplied by six}

Two issues must be repaired. First, a value jump in $L^2$ has no bounded
right inverse into $H^1$, whose trace lies in $H^{1/2}$; a patched
Laplacian residual needs value and first-derivative overlap bounds. Second,
at $\lambda>1$ a small full-space residual does not imply a discrete
eigenvalue, because a continuous-spectrum wave packet is a counterexample.
The construction below uses smooth local averaging and an explicit
cuspidal projection.

\section{Finite-index induction}

Put $\G=\mathrm{PSL}_2(\ZZ[i])$ and $\Hsub\subset\G$. Let
$C=\Hsub\backslash\G\simeq\mathbb P^1(\mathbb F_5)$, ordered as
\[
 (0:1),(1:0),(1:1),(1:2),(1:3),(1:4).
\]
For the right action on $C$, define the unitary permutation representation
\begin{equation}\label{eq:rho}
 (\rhoind(\gamma)v)_c=v_{c\cdot\gamma}.
\end{equation}
The computational arrays store
$\pi_\gamma(c)=c\cdot\gamma^{-1}$. Thus
\[
 F_c(P)=F_{\pi_\gamma(c)}(\gamma P)
 \quad\Longleftrightarrow\quad
 F(\gamma P)=\rhoind(\gamma)F(P).
\]

\begin{theorem}[Unitary induction]\label{thm:induction}
For right-coset representatives $g_c$, the map
\[
 (\mathcal Jf)_c(P)=f(g_cP)
\]
is a unitary isomorphism
\[
 L^2(\Hsub\backslash\HH)
 \simeq L^2(\G\backslash\HH;\CC^6_{\rhoind}).
\]
It intertwines the self-adjoint Laplacians and maps the scalar cuspidal
subspace onto the vector-valued cuspidal subspace.
\end{theorem}

\begin{proof}
Write $g_c\gamma=h g_{c\cdot\gamma}$ with $h\in\Hsub$ to obtain
\eqref{eq:rho}. If $F_1$ is a fundamental domain for $\G$, then
$\bigsqcup_c g_cF_1$ is, up to its faces, a fundamental domain for
$\Hsub$, so invariance of hyperbolic measure proves equality of norms.
Every $g_c$ is an isometry, proving the domain and Laplacian statements.
Finally, the double cosets $\Hsub\backslash\G/\G_\infty$ are the cusps of
$\Hsub$. Averaging components over each parabolic orbit gives exactly the
scalar constant term at that cusp, proving the cuspidal statement.
\end{proof}

\begin{proposition}[Two singular channels]\label{prop:channels}
For the exact repository permutations, both the translation image and the
full cusp-stabilizer image have orbits
\[
 \{0\},\qquad\{1,2,3,4,5\}.
\]
Hence the singular subspace is two-dimensional, with basis
\[
 u_\infty=e_0,\qquad
 u_0=5^{-1/2}(e_1+\cdots+e_5),
\]
and projector
\[
 P_{\rm sing}=e_0e_0^*
 +\frac15\begin{pmatrix}0&0\\0&\mathbf1_5\mathbf1_5^*\end{pmatrix}.
\]
The other four dimensions are regular parabolic character channels.
\end{proposition}

\begin{proof}
Invert the stored glue maps to get the forward actions and derive
$T_i$ from $T_iR=T_i\,R$. Exact permutation closure gives translation
image order five and full parabolic image order ten, with the stated
orbits. Fixed vectors of a permutation group are the vectors constant on
its orbits. The integer/rational verification is
\texttt{six\_copy\_DK\_algebra.py}.
\end{proof}

\begin{remark}[Why $T_iR$ is retained]
Exact closure shows that the stored $T_iR$ permutation already lies in
$\langle T_1,R\rangle$, so it does not enlarge the cusp-stabilizer image.
It is retained because it is an actual repository face transition and gives
an independent consistency check when deriving the $T_i$ action. As an
ambient check, the four stored permutations generate a transitive group of
order $60$ on the six cosets, as expected from
$\mathrm{PSL}_2(\mathbb F_5)$.
\end{remark}

\begin{remark}
There are two cusp channels, not six independent cusp tails. A factor
$\sqrt6$ enters only when a componentwise maximum is converted to the
Euclidean norm on $\CC^6$.
\end{remark}

\section{Smooth gluing from certified overlap defects}

Let $X=\G\backslash\HH$ and let $E_{\rhoind}\to X$ be the induced flat
bundle, with flat unitary connection $\nabla$. Fix a locally finite
orbifold uniformizer cover $(\widetilde U_j,G_j)$ with only finitely many
transition types: finitely many
compact-core charts and one periodic cusp chart type. Choose fixed lifts and
explicit unitary flat-bundle transitions $t_{kj}$ from the $j$-trivialization
to the $k$-trivialization on $U_j\cap U_k$. Let
$\psi_j\in C^2(X)$ be nonnegative, locally finite, subordinate to the
cover, and $\sum_j\psi_j=1$. Require that every $\psi_jF_j$, initially
defined on $U_j$, extends by zero to a $C^2$ global section. This holds, for
example, when $\supp\psi_j\Subset U_j$. The periodic cusp version may instead
be verified directly on one period and translated. The certificate must prove
local finiteness, coverage, cusp periodicity, the extension property, and
integrability of the derivative majorants.

Let $W_j$ be the raw finite Whittaker field on $\widetilde U_j$, and let
$\rho_j$ be the finite isotropy action on its bundle fiber. Define
\begin{equation}\label{eq:isotropyaverage}
 F_j(x)=\frac1{\abs{G_j}}\sum_{h\in G_j}\rho_j(h)^{-1}W_j(hx).
\end{equation}
Then $F_j(gx)=\rho_j(g)F_j(x)$, so $F_j$ is a genuine smooth local
orbibundle section. Since every $h$ is an isometry and the fiber action is
constant and unitary, isotropy averaging commutes with $\D$; in particular
$(\D-\lambda)F_j=0$ whenever the finite Whittaker terms satisfy their local
eigen-equation. At a regular chart $G_j$ is trivial. For a measurable choice
of active anchor $k=k(x)$, put
$F_j^{(k)}=t_{kj}F_j$ and
\begin{align}
 d_0(x)&=\max_{j:\psi_j(x)\ne0}
   \abs{F_j^{(k)}(x)-F_k(x)}_{\CC^6},\\
 d_1(x)&=\max_{j:\psi_j(x)\ne0}
   \abs{\nabla F_j^{(k)}(x)-\nabla F_k(x)}_{T^*X\otimes\CC^6},\\
 B_0(x)&=\sum_j\abs{\D\psi_j(x)},\qquad
 B_1(x)=2\sum_j\abs{\nabla\psi_j(x)}.
\end{align}

\begin{lemma}[Certified automorphization]\label{lem:patch}
The expression
\[
 U=\sum_j\psi_jF_j
\]
is an exactly equivariant global section. Let $F_*(x)=F_{k(x)}(x)$ in the
chosen active-anchor trivialization and
\[
 \tau=\norm{\sum_j\psi_j(\D-\lambda)F_j}_2,
\]
where the displayed sum is formed in the common fiber. Assume
$F_*\in L^2$, $d_0\in L^2$, and
$\tau+\norm{B_0d_0+B_1d_1}_2<\infty$. Then $U\in\Dom(\D)$ and
\begin{align}
 \norm{U-F_*}_2&\le\norm{d_0}_2,\label{eq:normloss}\\
 R:=\norm{(\D-\lambda)U}_2
 &\le\tau+\norm{B_0d_0+B_1d_1}_2.\label{eq:patchres}
\end{align}
In particular, certified uniform bounds $d_i\le\delta_i$ give
\begin{equation}\label{eq:patchsup}
 R\le\tau+b_0\delta_0+b_1\delta_1,\qquad
 b_i=\norm{B_i}_2.
\end{equation}
For the finite Whittaker field, $\tau=0$ termwise.
\end{lemma}

\begin{proof}
Each $F_j$ is a local section and each weighted term has a $C^2$ extension
by zero, so the locally finite partition sum is a global section. In the
$k(x)$ trivialization, convexity of the weights gives
\eqref{eq:normloss}. Apply the covariant Laplacian product rule in the same
trivialization. Terms containing the common anchor $F_k$ and $\nabla F_k$
cancel because $\sum_j\D\psi_j=0$ and
$\sum_j\nabla\psi_j=0$. Unitarity and flatness of $t_{kj}$ make the
remaining commutators bounded pointwise by $B_0d_0+B_1d_1$, proving
\eqref{eq:patchres} distributionally. The hypotheses and
\eqref{eq:normloss} give $U\in L^2$, while the right side of
\eqref{eq:patchres} gives the distributional $\D U\in L^2$. On the complete
finite-volume orbifold, the maximal $L^2$ domain of the flat unitary
Laplacian is its self-adjoint
domain~\cite[Remark~3.14]{BallmannPolymerakis}. Hence $U\in\Dom(\D)$.
\end{proof}

\begin{remark}[The current continuum number is not yet $\delta_0$]
The current computation takes a maximum over scalar components and four
generators on a rectangular core superset. Lemma~\ref{lem:patch} needs the
vector norm on every transition active in the certified cover, including
edge and vertex transition words. It also needs $d_1$ and the partition
constants $b_0,b_1$. At elliptic charts the candidates must first be averaged
as in \eqref{eq:isotropyaverage}; the resulting changes are included in the
same value and covariant-gradient defects. A scalar component maximum costs
at most $\sqrt6$ when converted to the vector norm.
\end{remark}

\section{Spectral enclosure and the continuous spectrum}

\begin{lemma}[Cuspidal quasimode]\label{lem:quasimode}
Let $P_{\rm cusp}$ be an orthogonal projection commuting with $\D$ and
having cuspidal range. If
\[
 \norm{P_{\rm cusp}U}_2\ge\mu>0,
\]
then some cuspidal eigenvalue satisfies
\begin{equation}\label{eq:eigbound}
 \abs{\lambda_j-\lambda}\le R/\mu.
\end{equation}
\end{lemma}

\begin{proof}
Expand $P_{\rm cusp}U$ in an orthonormal basis of the discrete cuspidal
spectrum. If every eigenvalue were farther than $R/\mu$, its squared
residual would exceed $R^2$. Projection contractivity and commutation with
$\D$ finish the proof.
\end{proof}

\subsection{Track B: the unconditional old, quarter-turn-odd sector}

The induced representation has the invariant line
\[
 V_{\rm old}=\CC u_{\rm old},\qquad
 u_{\rm old}=6^{-1/2}(1,1,1,1,1,1)^t.
\]
Its orthogonal projector $P_{\rm old}=u_{\rm old}u_{\rm old}^*$ commutes
with the representation and the Laplacian. Under
Theorem~\ref{thm:induction}, this is exactly the level-one oldspace.

Let $Q(z,y)=(iz,y)$. This isometry normalizes the Picard group and its
parabolic subgroup, fixes its cusp and height, and commutes with the
Laplacian. Although $Q$ has order four on $\HH$, its square is
\[
 Q^2(z,y)=(-z,y)=
 \begin{pmatrix}i&0\\0&-i\end{pmatrix}(z,y),
 \qquad
 \begin{pmatrix}i&0\\0&-i\end{pmatrix}\in\G.
\]
Consequently the scalar pullback $\mathcal Qf=f\circ Q$ satisfies
$\mathcal Q^2=I$. It is unitary and therefore self-adjoint, so
$(I-\mathcal Q)/2$ is an orthogonal projector with no simplicity hypothesis.
Let
$I_{\rm old}f=u_{\rm old}f$ be the isometric identification with the
six-copy oldspace and set
\[
 Q_{\rm old}=I_{\rm old}\mathcal QI_{\rm old}^*
 \quad\hbox{on }\Ran P_{\rm old}.
\]
Extend it by the identity on $\Ran P_{\rm old}^{\perp}$, calling the result
$\widehat Q$. Thus the oldspace reduces $\widehat Q$, and
$P_{Q,-}=(I-\widehat Q)/2$ satisfies
$P_{Q,-}^2=P_{Q,-}=P_{Q,-}^*$ and commutes with $P_{\rm old}$.

\begin{proposition}[Quarter-turn-odd oldspace is cuspidal]\label{prop:odd}
The range of
\[
 P_{\rm target}=P_{Q,-}P_{\rm old}
\]
is contained in the cuspidal subspace.
\end{proposition}

\begin{proof}
The level-one Picard quotient has one singular cusp channel. Its Eisenstein
family is induced from constant cusp data. The quarter-turn fixes that datum
and reindexes the defining sum over $\G_\infty\backslash\G$, so the
Eisenstein series is $Q$-even in its half-plane of convergence. Meromorphic
continuation makes the entire Eisenstein family and all of its residues
$Q$-even. Hence the $Q=-1$ level-one subspace is orthogonal to the full
continuous and residual noncuspidal spectrum and is cuspidal. A level-one
cusp form remains cuspidal at every subgroup cusp: each subgroup constant
term maps to the constant term at the unique level-one cusp. The oldspace
identification therefore transfers the conclusion to the six-copy local
system. Since $P_{Q,-}$ and $P_{\rm old}$ are commuting orthogonal projections,
$P_{\rm target}$ is itself the orthogonal projection required in
Lemma~\ref{lem:quasimode}; by construction it also commutes with $\D$.
\end{proof}

\begin{corollary}[Track-B six-copy D$(K)$]\label{cor:target}
If interval arithmetic proves
\[
 \norm{P_{\rm target}U}_2\ge\mu_{\rm target}>0,
\]
then a cusp form on $\Hsub\backslash\HH$ has an eigenvalue satisfying
\begin{equation}\label{eq:targetbound}
 \abs{\lambda_j-\lambda}
 \le\frac{\tau+b_0\delta_0+b_1\delta_1}{\mu_{\rm target}}.
\end{equation}
No Hecke growth hypothesis or soft continuous-spectrum cutoff enters this
bound.
\end{corollary}

\begin{proposition}[Certified Track-B cusp witness]\label{prop:trackBwitness}
Let $\Omega_B\subset X$ be a measurable cusp slab and let $W_B$ be a
certified local section there in the range of $P_{\rm target}$. If interval
arithmetic proves
\[
 \norm{W_B}_{L^2(\Omega_B)}\ge m_B,
 \qquad
 \norm{(P_{\rm target}U)|_{\Omega_B}-W_B}_{L^2(\Omega_B)}
 \le\varepsilon_B<m_B,
\]
then Corollary~\ref{cor:target} is valid with
\begin{equation}\label{eq:trackBmass}
 \mu_{\rm target}=m_B-\varepsilon_B>0.
\end{equation}
\end{proposition}

\begin{proof}
Restriction can only decrease the global norm. The reverse triangle
inequality on $\Omega_B$ gives
\[
 \norm{P_{\rm target}U}_{L^2(X)}
 \ge \norm{P_{\rm target}U}_{L^2(\Omega_B)}
 \ge m_B-\varepsilon_B.
\]
\end{proof}

\begin{corollary}[Track-B plateau closes the mass input]
\label{cor:trackBplateau}
Let $\mathcal H_1$ be the level-one cusp chart $y>1$. Define the exact
$C^2$ cutoff
\[
 \chi_B(y)=
 \begin{cases}
 0,&y\le1.01,\\
 10t^3-15t^4+6t^5,&1.01<y<1.20,\\
 1,&y\ge1.20,
 \end{cases}
 \qquad t=\frac{y-1.01}{1.20-1.01}.
\]
Let $\{\phi_j\}$ be a normalized subordinate family for all other charts,
with $\sum_j\phi_j=1$ wherever $1-\chi_B\ne0$. Use $W_B$ with weight
$\chi_B$ and every other local candidate with weight
$(1-\chi_B)\phi_j$. These weights still sum to one, and the glued section
satisfies
$U=W_B=P_{\rm target}U$ on
\[
 \Omega_B=[-1/2,1/2]\times[0,1/2]\times[1.20,1.45].
\]
Consequently $\varepsilon_B=0$ in
Proposition~\ref{prop:trackBwitness}. For the fixed recorded trial, the Arb
certificate gives the admissible mass input
\begin{equation}\label{eq:trackBnumericmass}
 \mu_{\rm target}>0.20380681000849094.
\end{equation}
\end{corollary}

\begin{proof}
The horoball $y>1$ embeds modulo the Picard parabolic subgroup. Indeed, for
$\gamma=\left(\begin{smallmatrix}a&b\\c&d\end{smallmatrix}\right)$ with
$c\ne0$,
\[
 y(\gamma P)=\frac{y}{\abs{cz+d}^2+\abs c^2y^2}
 \le\frac1{\abs c^2y}<1
 \quad(y>1),
\]
because $\abs c\ge1$. Thus the closure of $\supp\chi_B$ remains inside a
valid cusp chart. The quintic and its first two derivatives match the
adjacent constants, so extension by zero is $C^2$. The integer Fourier
lattice, $-z$ average, and oldspace vector make $W_B$ a genuine cusp local
section; its construction also gives $P_{\rm target}W_B=W_B$. On the plateau
all other weights vanish, hence $U=W_B$ there. Since $Q$ preserves height, it
also preserves the plateau, so applying $P_{\rm target}$ does not sample
outside this identity region. The numerical lower bound is the 192-bit,
512-height-segment Arb result recorded by
\texttt{track\_b\_projected\_mass\_arb.py}.
\end{proof}

\begin{remark}[Rotation odd is sufficient, but is not a class label]
In Then's four-sign notation, $P_{Q,-}$ contains the $C$ and $H$ sectors;
the reflection sign distinguishes them. Corollary~\ref{cor:target} only
asserts a nearby cuspidal eigenvalue and does not label its symmetry class.
If a certified class-$C$ identification is later required, one may append
the commuting reflection-even projector. It is unnecessary for removing the
continuous spectrum.
\end{remark}

Then's table places the published
$r=6.6221193402528\ldots$ Picard mode in class $C$, characterized by
$f(iz,y)=-f(z,y)$ and $f(-\overline z,y)=f(z,y)$. The corrected six-copy
coefficients are numerically a level-one lift.
These facts motivate Corollary~\ref{cor:target}, but do not replace the
certified projected-mass lower bound.

For the fixed $M=8$ coefficient file, the non-certifying coefficient
diagnostic gives relative defects $1.37\cdot10^{-3}$ for
$a_{i\beta}=-a_\beta$ and $1.28\cdot10^{-3}$ for
$a_{\overline\beta}=a_\beta$. The zero-cusp off-level-one fraction is
$5.20\cdot10^{-4}$ and the matched-cusp coefficient mismatch is
$6.08\cdot10^{-4}$. On a quarter-turn-invariant horosphere sample, the
oldspace and old-class-$C$ projectors each retain $0.9999998893$ of the
sampled mass. These diagnostics agree with, but are not used by, the
certified Track-B mass bound \eqref{eq:trackBnumericmass}.

\subsection{General two-cusp route}

\begin{assumption}[HQ: normalized Hecke annihilation]\label{ass:HQ}
For a Gaussian prime ideal $\mathfrak q\nmid(2+i)$ there is a bounded
self-adjoint Hecke operator $T_{\mathfrak q}$ commuting with $\D$ and an
explicit real $C^1$ function $h_{\mathfrak q}$ on $[0,\infty)$ such that
\[
 T_{\mathfrak q}=h_{\mathfrak q}(\D)
\quad\text{on the full Eisenstein and residual spectrum.}
\]
Certified constants satisfy
\[
 \norm{h_{\mathfrak q}}_\infty\le M_h,\qquad
 \norm{h'_{\mathfrak q}}_\infty\le L_h,\qquad
 \norm{T_{\mathfrak q}}\le M_T.
\]
The function $h_{\mathfrak q}$ must be derived from the chosen Bianchi
double-coset normalization; it must not be imported by analogy from
$\mathrm{PSL}_2(\ZZ)$.
\end{assumption}

Under HQ,
\[
 \Diamond_{\mathfrak q}=h_{\mathfrak q}(\D)-T_{\mathfrak q}
\]
is self-adjoint, commutes with $\D$, and vanishes
on the full continuous and residual subspace. Its range is therefore
orthogonal to that subspace, hence cuspidal, and its norm is at most
$M_h+M_T$. For a certified real number
$a_{\mathfrak q}$ and Hecke defect
\[
 H_{\mathfrak q}=\norm{(T_{\mathfrak q}-a_{\mathfrak q})U}_2,
\]
functional calculus gives
\begin{equation}\label{eq:heckemass}
 \norm{\Diamond_{\mathfrak q}U}_2
 \ge
 \abs{h_{\mathfrak q}(\lambda)-a_{\mathfrak q}}\norm{U}_2
 -L_hR-H_{\mathfrak q}.
\end{equation}
Moreover, commutation with $\D$ gives
\[
 \norm{(\D-\lambda)\Diamond_{\mathfrak q}U}_2
 \le \norm{\Diamond_{\mathfrak q}}R
 \le (M_h+M_T)R.
\]

\begin{corollary}[General conditional six-copy D$(K)$]\label{cor:general}
Assume HQ and suppose the right side of \eqref{eq:heckemass} has certified
positive lower bound $\mu_{\mathfrak q}$. Then
\[
 \abs{\lambda_j-\lambda}
 \le
 \frac{(M_h+M_T)(\tau+b_0\delta_0+b_1\delta_1)}
 {\mu_{\mathfrak q}}
\]
for a cusp-form eigenvalue of $\Hsub$.
\end{corollary}

\section{Proof-producing certificate interface}

The target branch requires the following interval-certified inputs:
\begin{enumerate}[leftmargin=*]
\item exact transition permutations, representation convention, finite
isotropy groups and their fiber actions;
\item a finite-type $C^2$ chart partition, coverage and local-finiteness
proofs, cusp periodicity, and bounds $b_0,b_1$;
\item vector overlap bounds $\delta_0,\delta_1$ for every active transition,
after the exact stabilizer averages \eqref{eq:isotropyaverage};
\item for Track B, the certified plateau witness of
Corollary~\ref{cor:trackBplateau}; its projected-mass input is closed with
$\varepsilon_B=0$;
\item proofs, in the correctly scaled coordinates of both subgroup cusps,
that every local expansion has zero constant Fourier coefficient and
exponential $L^2$ decay after flat-bundle chart transport;
\item periodicity/local finiteness of the cusp partition and an $L^2$
majorant for $B_0d_0+B_1d_1$. The finite Whittaker identity already gives
$\tau=0$ termwise.
\end{enumerate}
The general branch additionally requires HQ, a Hecke defect, and the
non-resonance lower bound in \eqref{eq:heckemass}.

The smallest next implementation is therefore not another SVD. It is an
interval-certified extension of the fixed cusp cutoff to a $C^2$ partition
on a thickened Humbert cell, enumeration of its finite transition set, and
interval automatic differentiation of the Whittaker field through first
order. The projected norm is already certified by
Corollary~\ref{cor:trackBplateau}.

\begin{thebibliography}{9}
\bibitem{BSV}
A.~R. Booker, A. Str\"ombergsson, and A. Venkatesh,
\emph{Effective computation of Maass cusp forms},
Int. Math. Res. Not. 2006, Art. ID 71281.

\bibitem{BallmannPolymerakis}
W.~Ballmann and P.~Polymerakis,
\emph{On the differential form spectrum of geometrically finite orbifolds},
MPIM Preprint 2020 (62), arXiv:2011.13304.

\bibitem{Friedman}
J.~S. Friedman,
\emph{The Selberg trace formula and Selberg zeta-function for cofinite
Kleinian groups with finite-dimensional unitary representations},
Ph.D. thesis, 2005, Chapter~3.

\bibitem{FriedmanArtin}
J.~S. Friedman,
\emph{Analogues of the Artin factorization formula for the automorphic
scattering matrix and Selberg zeta-function associated to a Kleinian group},
Proc. Amer. Math. Soc. 2008.

\bibitem{Then}
H.~Then, \emph{Arithmetic quantum chaos of Maass waveforms}, Picard-group
tables and symmetry classes.
\end{thebibliography}
\end{document}
